\section{排序、选择和区间章正文案例推导补强}

这一节补齐本章正文出现过的 LeetCode 案例推导。阅读顺序统一按“题目说明、样例入口、暴力基线、暴力过程、重复工作、优化条件、相邻状态、状态复用、指针和字段、代码落点、正确性、边界实验、复杂度变化、迁移追问”展开；目的不是新增题量，而是把正文中的案例都讲成可复述、可证明、可迁移的解题链。

\subsection{LeetCode 75 颜色分类}

\textbf{题目说明。} LeetCode 75 颜色分类 的题面入口要先还原成具体任务：三指针维护零区、一号区、未知区和二区。

\textbf{样例入口。} 拿 [2, 0, 1] 手算，mid 指向 2 时和右边界交换，换回来的元素未知，所以 mid 不能前进；再看到 1 才前进。

\textbf{暴力基线。} 暴力可以计数 0、1、2 的个数后重写数组，或者排序数组。

\textbf{暴力过程。} 以 [2, 0, 1] 为例，计数法先数出每种颜色，再写 0、1、2。原地一趟要在扫描时把 0 放左边、2 放右边。

\textbf{重复工作。} 题面模型：三指针维护零区、一号区、未知区和二区。暴力版的慢点要落到本题：浪费在两趟计数或排序。颜色只有三类，可以在一趟扫描中维护四个区域。后面的优化只消掉这类重复工作，不能改变答案集合。

\textbf{优化条件。} 本题的关键观察是：遇到二时只移动右边界，不移动扫描指针，因为换回来的元素未知。这句话必须能翻译成可维护状态；如果题目条件改变后观察不再成立，就不能继续套原来的写法。

\textbf{相邻状态。} 规律是 [0, low) 是 0， [low, mid) 是 1， (high, n) 是 2，未知区只在 mid 到 high 之间收缩。把这个特例继续拆开看：先标出已经确认的前缀或后缀，再标出未知区；能被丢掉的候选，必须是以后再也不会改变已确认区域语义的候选。这样才算从例子推出数组不变量，而不是只记一次操作。

\textbf{状态复用。} 维护 low、mid、high：[0, low) 是 0，[low, mid) 是 1，[mid, high] 是未知，(high, n) 是 2。手算时只改被新输入影响的那一小部分，而不是回头重算完整候选。

\textbf{指针和字段。} 状态字段要能说清：mid 是当前检查位置，low 是下一个 0 的位置，high 是下一个 2 的位置。手算时按这个协议跟踪：[2, 0, 1] 中 mid 指向 2，和 high 交换后 high--，mid 不动，因为换回来的 1 还没检查；再处理 1，mid++。

\textbf{代码落点。} 关键是遇到 0 时 swap(nums[low++], nums[mid++])；遇到 2 时 swap(nums[mid], nums[high--])，mid 不能前进。关键代码行必须能逐句翻译成“旧状态怎样变成新状态”，否则就只是模板。

\textbf{正确性。} 每次操作都把一个元素放入 0 区或 2 区，或确认它属于 1 区；未知区严格缩小，不变量保持。

\textbf{边界实验。} 先用全零、全二、已经有序、逆序、交换后漏检查做反例检查。实验时把全零、全二、已经有序、逆序、交换后漏检查列成表，再打印关键下标和有效区间。若某个元素被写入后又被未来步骤破坏，说明区间不变量没有成立。

\textbf{复杂度变化。} 排序 O(n log n)，计数 O(n) 两趟；荷兰国旗算法 O(n) 时间、O(1) 空间。

\textbf{迁移追问。} 如果题目改成“这是荷兰国旗问题，能推广到三类分区的原地维护”，先判断原状态是否仍然覆盖新答案集合，再决定是微调边界、改 value 语义，还是换算法。

\subsection{LeetCode 56 合并区间}

\textbf{题目说明。} LeetCode 56 合并区间 的题面入口要先还原成具体任务：排序后可能与当前区间相交的只会是结果尾部。

\textbf{样例入口。} 拿 [[1, 3], [2, 6], [8, 10]] 手算，排序后 [2, 6] 只可能和结果尾部 [1, 3] 相交，合并成 [1, 6]；[8, 10] 的起点大于尾部右端，另开新区间。

\textbf{暴力基线。} 暴力可以不断挑一个区间，和所有其他区间比较是否相交，能合并就合并，直到没有变化。

\textbf{暴力过程。} 以 [[1, 3], [2, 6], [8, 10]] 为例，[1, 3] 和 [2, 6] 相交，应合成 [1, 6]；[8, 10] 与它不相交，单独保留。

\textbf{重复工作。} 题面模型：排序后可能与当前区间相交的只会是结果尾部。暴力版的慢点要落到本题：浪费在反复两两比较。按左端排序后，能和当前结果尾部相交的区间会连续出现，不需要再回头看更早区间。后面的优化只消掉这类重复工作，不能改变答案集合。

\textbf{优化条件。} 本题的关键观察是：按起点排序，扫描时维护结果最后区间的右端点。这句话必须能翻译成可维护状态；如果题目条件改变后观察不再成立，就不能继续套原来的写法。

\textbf{相邻状态。} 规律是排序后当前区间只需要看结果尾部，端点相接是否合并要按闭区间语义决定。把这个特例继续拆开看：排序以后，真正还能影响当前区间的候选必须贴近当前端点、结果尾部或资源堆顶。某个区间一旦被覆盖、合并或弹出，就要说明它为什么不会和后续区间重新产生更优关系。

\textbf{状态复用。} 先按 start 升序排序。结果数组保持已合并且不重叠；当前区间只和结果最后一个区间比较。手算时只改被新输入影响的那一小部分，而不是回头重算完整候选。

\textbf{指针和字段。} 状态字段要能说清：last 是结果尾区间，interval 是当前扫描区间；若 interval.start<=last.end 就合并，否则追加新区间。手算时按这个协议跟踪：排序后先放 [1, 3]；看到 [2, 6]，2<=3，更新尾部右端为 6；看到 [8, 10]，8>6，追加。

\textbf{代码落点。} 关键是明确闭区间语义，端点相接是否合并由题目决定。本题闭区间重叠时用 start<=last.end。关键代码行必须能逐句翻译成“旧状态怎样变成新状态”，否则就只是模板。

\textbf{正确性。} 排序保证当前区间之前的所有区间左端不大于它；若它不和结果尾相交，就不可能和更早且右端不超过尾部的区间相交。

\textbf{边界实验。} 先用端点相接、完全覆盖、无重叠、空列表、输入未排序做反例检查。实验时覆盖端点相接、完全覆盖、无重叠、空列表、输入未排序，画出排序后的区间序列和当前结果尾部。动态版本要专门测前驱、后继和端点相接。

\textbf{复杂度变化。} 暴力反复合并最坏 O(\ensuremath{n^{2}})，排序扫描 O(n log n) 时间，输出空间 O(n)。

\textbf{迁移追问。} 如果题目改成“若区间是半开区间，端点相等不一定合并”，先判断原状态是否仍然覆盖新答案集合，再决定是微调边界、改 value 语义，还是换算法。

\subsection{LeetCode 57 插入区间}

\textbf{题目说明。} LeetCode 57 插入区间 的题面入口要先还原成具体任务：新区间只会影响与它相交的一段连续区间。

\textbf{样例入口。} 拿 intervals=[[1, 3], [6, 9]]、new=[2, 5] 手算，左侧完全小于新区间的先放入；[1, 3] 与 [2, 5] 相交，合成 [1, 5]；[6, 9] 在右侧直接追加。

\textbf{暴力基线。} 慢版本对新区间和所有旧区间两两比较，手工判断相交、覆盖或冲突。它能验证端点语义，但排序后很多远离当前边界的区间无需再看。放到本题就是：先按“新区间只会影响与它相交的一段连续区间”完整试慢版本；样例里已经能看到“规律是新区间影响的只是一段连续重叠区间，先左、再合并、再右，原列表为空或新区间覆盖全部也服从同一流程”，所以优化前必须先能说清慢版本枚举了哪些候选。

\textbf{暴力过程。} 拿 intervals=[[1, 3], [6, 9]]、new=[2, 5] 手算，左侧完全小于新区间的先放入；[1, 3] 与 [2, 5] 相交，合成 [1, 5]；[6, 9] 在右侧直接追加。

\textbf{重复工作。} 题面模型：新区间只会影响与它相交的一段连续区间。暴力版的慢点要落到本题：慢版本会围绕“新区间只会影响与它相交的一段连续区间”做两两区间比较。样例规律是“规律是新区间影响的只是一段连续重叠区间，先左、再合并、再右，原列表为空或新区间覆盖全部也服从同一流程”，说明排序后只有贴近当前端点、结果尾部或资源堆顶的区间还会影响答案。后面的优化只消掉这类重复工作，不能改变答案集合。

\textbf{优化条件。} 本题的关键观察是：先追加左侧完全不相交区间，再合并重叠段，最后追加右侧。这句话必须能翻译成可维护状态；如果题目条件改变后观察不再成立，就不能继续套原来的写法。

\textbf{相邻状态。} 规律是新区间影响的只是一段连续重叠区间，先左、再合并、再右，原列表为空或新区间覆盖全部也服从同一流程。把这个特例继续拆开看：排序以后，真正还能影响当前区间的候选必须贴近当前端点、结果尾部或资源堆顶。某个区间一旦被覆盖、合并或弹出，就要说明它为什么不会和后续区间重新产生更优关系。

\textbf{状态复用。} 把两两比较压成排序后的相邻关系、当前边界或动态有序结构；远离当前边界的区间要被安全归档。本题落点是：规律是新区间影响的只是一段连续重叠区间，先左、再合并、再右，原列表为空或新区间覆盖全部也服从同一流程。手算时只改被新输入影响的那一小部分，而不是回头重算完整候选。

\textbf{指针和字段。} 状态字段要能说清：排序后的区间、当前结果尾、扫描右端或资源堆、端点比较规则；动态版本还要保存前驱后继。手算时按这个协议跟踪：拿 intervals=[[1, 3], [6, 9]]、new=[2, 5] 手算，左侧完全小于新区间的先放入；[1, 3] 与 [2, 5] 相交，合成 [1, 5]；[6, 9] 在右侧直接追加。然后按协议复查：手算时先排序，再逐个区间写当前结果尾部、当前右端或资源堆。每个区间离开视野时都要说明它不会再影响后续。

\textbf{代码落点。} 排序比较器满足严格弱序；端点相等的处理和题意一致；扫描时只和当前结果尾或当前资源集合交互。关键代码行必须能逐句翻译成“旧状态怎样变成新状态”，否则就只是模板。

\textbf{正确性。} 证明从排序后的邻接关系开始：已经合并、选择或丢弃的区间不会被更后面的区间重新影响。

\textbf{边界实验。} 先用新区间在最前、最后、覆盖全部、刚好相邻、原列表为空做反例检查。实验时覆盖新区间在最前、最后、覆盖全部、刚好相邻、原列表为空，画出排序后的区间序列和当前结果尾部。动态版本要专门测前驱、后继和端点相接。

\textbf{复杂度变化。} 暴力两两比较区间常见为平方级；排序后扫描通常由排序成本主导，动态有序结构则把每次前驱后继查询控制在对数级。

\textbf{迁移追问。} 如果题目改成“若输入不是有序且不重叠，必须先排序或换模型”，先判断原状态是否仍然覆盖新答案集合，再决定是微调边界、改 value 语义，还是换算法。

\subsection{LeetCode 435 无重叠区间}

\textbf{题目说明。} LeetCode 435 无重叠区间 的题面入口要先还原成具体任务：保留最多不重叠区间等价于删除最少区间。

\textbf{样例入口。} 拿 [[1, 2], [2, 3], [3, 4], [1, 3]] 手算，选 [1, 2] 后还能接 [2, 3] 和 [3, 4]；若先选 [1, 3]，会挡住更多区间。

\textbf{暴力基线。} 暴力可以枚举保留哪些区间，再检查它们是否互不重叠，最后用总数减最大保留数量。

\textbf{暴力过程。} 以 [[1, 2], [2, 3], [3, 4], [1, 3]] 为例，选 [1, 2] 后还能接 [2, 3] 和 [3, 4]；若先选 [1, 3] 会挡住更多后续区间。

\textbf{重复工作。} 题面模型：保留最多不重叠区间等价于删除最少区间。暴力版的慢点要落到本题：浪费在枚举保留集合。区间选择里结束越早，给后面留下的空间越大，可以作为安全局部选择。后面的优化只消掉这类重复工作，不能改变答案集合。

\textbf{优化条件。} 本题的关键观察是：按终点升序选择，结束越早给后续留下空间越多。这句话必须能翻译成可维护状态；如果题目条件改变后观察不再成立，就不能继续套原来的写法。

\textbf{相邻状态。} 规律是按右端点升序选择，结束越早给未来留下的空间越多，删除数等于总数减保留数。把这个特例继续拆开看：先构造一个非贪心选择，再比较两者留下的资源、右边界或剩余时间。只有能把非贪心第一步交换成贪心第一步且不变差，局部选择才是规律。

\textbf{状态复用。} 按右端点升序排序。维护 last\_end；若当前区间 start>=last\_end，就保留它并更新 last\_end，否则删除它。手算时只改被新输入影响的那一小部分，而不是回头重算完整候选。

\textbf{指针和字段。} 状态字段要能说清：last\_end 是已保留区间的最后右端，kept 是保留数量，removed 是删除数量。手算时按这个协议跟踪：排序后先保留 [1, 2]；[1, 3] 与它重叠删除；[2, 3] 和 [3, 4] 都能接上。

\textbf{代码落点。} 关键是端点相接是否算重叠。本题中 [1, 2] 和 [2, 3] 不重叠，所以条件是 start>=last\_end。关键代码行必须能逐句翻译成“旧状态怎样变成新状态”，否则就只是模板。

\textbf{正确性。} 任意最优解的第一个区间都可替换为结束最早的区间，后续可用空间不变小；重复此过程得到贪心解。

\textbf{边界实验。} 先用端点相接是否重叠、完全覆盖、相同终点、空列表做反例检查。实验时除了端点相接是否重叠、完全覆盖、相同终点、空列表，还要故意替换成一个看似合理但错误的局部规则，构造反例。能解释错误贪心为何失败，才算理解正确贪心。

\textbf{复杂度变化。} 暴力指数级；排序 O(n log n)，扫描 O(n)，空间取决于排序。

\textbf{迁移追问。} 如果题目改成“用交换论证说明最早结束的选择不劣于其他第一选择”，先判断原状态是否仍然覆盖新答案集合，再决定是微调边界、改 value 语义，还是换算法。

\subsection{LeetCode 215 数组中的第 K 个最大元素}

\textbf{题目说明。} LeetCode 215 数组中的第 K 个最大元素 的题面入口要先还原成具体任务：只关心第 K 大，不需要完整排序。

\textbf{样例入口。} 拿 [3, 2, 1, 5, 6, 4]、k=2 手算，固定大小为 2 的小顶堆最终保存最大的两个数，堆顶就是第 2 大。

\textbf{暴力基线。} 暴力排序整个数组，然后返回降序第 k 个或升序第 n-k 个。

\textbf{暴力过程。} 以 [3, 2, 1, 5, 6, 4]、k=2 为例，完整排序后第 2 大是 5；但为了一个第 K 大，不需要知道所有元素的完整顺序。

\textbf{重复工作。} 题面模型：只关心第 K 大，不需要完整排序。暴力版的慢点要落到本题：浪费在完整排序。我们只需要维护当前最大的 K 个元素，其中最小的那个就是第 K 大边界。后面的优化只消掉这类重复工作，不能改变答案集合。

\textbf{优化条件。} 本题的关键观察是：大小为 K 的小顶堆保存当前前 K 大边界；快速选择可做到平均线性。这句话必须能翻译成可维护状态；如果题目条件改变后观察不再成立，就不能继续套原来的写法。

\textbf{相邻状态。} 规律是堆里只保留当前前 k 大候选，比堆顶小的数不可能进入答案。把这个特例继续拆开看：堆顶必须正好代表下一步最需要处理的候选；若堆里可能有过期对象，读取堆顶前要先清理。堆只解决动态极值，不自动证明被弹出或被丢弃的元素无用。

\textbf{状态复用。} 大小为 K 的小顶堆保存当前前 K 大。新数大于堆顶时，弹出堆顶并加入新数；小于等于堆顶则不可能进入前 K。手算时只改被新输入影响的那一小部分，而不是回头重算完整候选。

\textbf{指针和字段。} 状态字段要能说清：heap.top 是当前前 K 大中的最小值，size 保持不超过 K，num 是当前扫描元素。手算时按这个协议跟踪：扫描样例时，堆最终保存 5 和 6，堆顶 5 就是第 2 大。

\textbf{代码落点。} 关键是使用小顶堆，C++ priority\_queue 默认大顶堆，需要 greater<int>。若选择快速选择，要处理随机 pivot 和重复值。关键代码行必须能逐句翻译成“旧状态怎样变成新状态”，否则就只是模板。

\textbf{正确性。} 堆中始终保存已扫描元素的前 K 大；若新元素不超过堆顶，它最多排在第 K 大之后，可以安全丢弃。

\textbf{边界实验。} 先用K 为一、K 为 n、重复值、随机 pivot做反例检查。实验时覆盖K 为一、K 为 n、重复值、随机 pivot，记录堆顶是否过期、比较器是否让正确候选在堆顶、K 或窗口大小为边界值时是否稳定。

\textbf{复杂度变化。} 排序 O(n log n)，小顶堆 O(n log k) 时间、O(k) 空间；快速选择平均 O(n)。

\textbf{迁移追问。} 如果题目改成“若数据流持续到来，堆方案比一次性快速选择更自然”，先判断原状态是否仍然覆盖新答案集合，再决定是微调边界、改 value 语义，还是换算法。

\subsection{LeetCode 253 会议室 II}

\textbf{题目说明。} LeetCode 253 会议室 II 的题面入口要先还原成具体任务：同时进行的会议数量决定最少房间数。

\textbf{样例入口。} 拿 [[0, 30], [5, 10], [15, 20]] 手算，0-30 占一间，5-10 来时最早结束仍是 30，所以需要第二间；15-20 来时 10 已结束，可复用。

\textbf{暴力基线。} 慢版本对新区间和所有旧区间两两比较，手工判断相交、覆盖或冲突。它能验证端点语义，但排序后很多远离当前边界的区间无需再看。放到本题就是：先按“同时进行的会议数量决定最少房间数”完整试慢版本；样例里已经能看到“规律是小顶堆保存当前占用会议的最早结束时间，堆大小就是房间数”，所以优化前必须先能说清慢版本枚举了哪些候选。

\textbf{暴力过程。} 拿 [[0, 30], [5, 10], [15, 20]] 手算，0-30 占一间，5-10 来时最早结束仍是 30，所以需要第二间；15-20 来时 10 已结束，可复用。

\textbf{重复工作。} 题面模型：同时进行的会议数量决定最少房间数。暴力版的慢点要落到本题：慢版本会围绕“同时进行的会议数量决定最少房间数”做两两区间比较。样例规律是“规律是小顶堆保存当前占用会议的最早结束时间，堆大小就是房间数”，说明排序后只有贴近当前端点、结果尾部或资源堆顶的区间还会影响答案。后面的优化只消掉这类重复工作，不能改变答案集合。

\textbf{优化条件。} 本题的关键观察是：按开始时间扫描，用小顶堆维护当前会议的最早结束时间。这句话必须能翻译成可维护状态；如果题目条件改变后观察不再成立，就不能继续套原来的写法。

\textbf{相邻状态。} 规律是小顶堆保存当前占用会议的最早结束时间，堆大小就是房间数。把这个特例继续拆开看：排序以后，真正还能影响当前区间的候选必须贴近当前端点、结果尾部或资源堆顶。某个区间一旦被覆盖、合并或弹出，就要说明它为什么不会和后续区间重新产生更优关系。

\textbf{状态复用。} 把两两比较压成排序后的相邻关系、当前边界或动态有序结构；远离当前边界的区间要被安全归档。本题落点是：规律是小顶堆保存当前占用会议的最早结束时间，堆大小就是房间数。手算时只改被新输入影响的那一小部分，而不是回头重算完整候选。

\textbf{指针和字段。} 状态字段要能说清：排序后的区间、当前结果尾、扫描右端或资源堆、端点比较规则；动态版本还要保存前驱后继。手算时按这个协议跟踪：拿 [[0, 30], [5, 10], [15, 20]] 手算，0-30 占一间，5-10 来时最早结束仍是 30，所以需要第二间；15-20 来时 10 已结束，可复用。然后按协议复查：手算时先排序，再逐个区间写当前结果尾部、当前右端或资源堆。每个区间离开视野时都要说明它不会再影响后续。

\textbf{代码落点。} 排序比较器满足严格弱序；端点相等的处理和题意一致；扫描时只和当前结果尾或当前资源集合交互。关键代码行必须能逐句翻译成“旧状态怎样变成新状态”，否则就只是模板。

\textbf{正确性。} 证明从排序后的邻接关系开始：已经合并、选择或丢弃的区间不会被更后面的区间重新影响。

\textbf{边界实验。} 先用会议端点相接、多个同一时间开始、空列表、堆顶释放条件做反例检查。实验时覆盖会议端点相接、多个同一时间开始、空列表、堆顶释放条件，画出排序后的区间序列和当前结果尾部。动态版本要专门测前驱、后继和端点相接。

\textbf{复杂度变化。} 暴力两两比较区间常见为平方级；排序后扫描通常由排序成本主导，动态有序结构则把每次前驱后继查询控制在对数级。

\textbf{迁移追问。} 如果题目改成“也可以把开始和结束分别排序，用双指针统计并发数”，先判断原状态是否仍然覆盖新答案集合，再决定是微调边界、改 value 语义，还是换算法。

\subsection{LeetCode 1109 航班预订统计}

\textbf{题目说明。} LeetCode 1109 航班预订统计 的题面入口要先还原成具体任务：每条预订给一个连续航班区间增加座位数。

\textbf{样例入口。} 拿 bookings=[[1, 2, 10], [2, 3, 20]]、n=3 手算，第 1 班加 10，第 2 班加 30，第 3 班加 20。

\textbf{暴力基线。} 慢版本对新区间和所有旧区间两两比较，手工判断相交、覆盖或冲突。它能验证端点语义，但排序后很多远离当前边界的区间无需再看。放到本题就是：先按“每条预订给一个连续航班区间增加座位数”完整试慢版本；样例里已经能看到“规律是区间加法用差分：起点加，终点后一位减，最后前缀还原”，所以优化前必须先能说清慢版本枚举了哪些候选。

\textbf{暴力过程。} 拿 bookings=[[1, 2, 10], [2, 3, 20]]、n=3 手算，第 1 班加 10，第 2 班加 30，第 3 班加 20。

\textbf{重复工作。} 题面模型：每条预订给一个连续航班区间增加座位数。暴力版的慢点要落到本题：慢版本会围绕“每条预订给一个连续航班区间增加座位数”做两两区间比较。样例规律是“规律是区间加法用差分：起点加，终点后一位减，最后前缀还原”，说明排序后只有贴近当前端点、结果尾部或资源堆顶的区间还会影响答案。后面的优化只消掉这类重复工作，不能改变答案集合。

\textbf{优化条件。} 本题的关键观察是：差分数组在区间起点加人数，终点后一位减人数，最后前缀还原每航班总数。这句话必须能翻译成可维护状态；如果题目条件改变后观察不再成立，就不能继续套原来的写法。

\textbf{相邻状态。} 规律是区间加法用差分：起点加，终点后一位减，最后前缀还原。把这个特例继续拆开看：排序以后，真正还能影响当前区间的候选必须贴近当前端点、结果尾部或资源堆顶。某个区间一旦被覆盖、合并或弹出，就要说明它为什么不会和后续区间重新产生更优关系。

\textbf{状态复用。} 把两两比较压成排序后的相邻关系、当前边界或动态有序结构；远离当前边界的区间要被安全归档。本题落点是：规律是区间加法用差分：起点加，终点后一位减，最后前缀还原。手算时只改被新输入影响的那一小部分，而不是回头重算完整候选。

\textbf{指针和字段。} 状态字段要能说清：排序后的区间、当前结果尾、扫描右端或资源堆、端点比较规则；动态版本还要保存前驱后继。手算时按这个协议跟踪：拿 bookings=[[1, 2, 10], [2, 3, 20]]、n=3 手算，第 1 班加 10，第 2 班加 30，第 3 班加 20。然后按协议复查：手算时先排序，再逐个区间写当前结果尾部、当前右端或资源堆。每个区间离开视野时都要说明它不会再影响后续。

\textbf{代码落点。} 排序比较器满足严格弱序；端点相等的处理和题意一致；扫描时只和当前结果尾或当前资源集合交互。关键代码行必须能逐句翻译成“旧状态怎样变成新状态”，否则就只是模板。

\textbf{正确性。} 证明从排序后的邻接关系开始：已经合并、选择或丢弃的区间不会被更后面的区间重新影响。

\textbf{边界实验。} 先用区间到最后一个航班、单点区间、多条重叠预订、下标从一开始做反例检查。实验时覆盖区间到最后一个航班、单点区间、多条重叠预订、下标从一开始，画出排序后的区间序列和当前结果尾部。动态版本要专门测前驱、后继和端点相接。

\textbf{复杂度变化。} 暴力两两比较区间常见为平方级；排序后扫描通常由排序成本主导，动态有序结构则把每次前驱后继查询控制在对数级。

\textbf{迁移追问。} 如果题目改成“这是静态区间加法，不需要线段树”，先判断原状态是否仍然覆盖新答案集合，再决定是微调边界、改 value 语义，还是换算法。

\subsection{LeetCode 347 前 K 个高频元素}

\textbf{题目说明。} LeetCode 347 前 K 个高频元素 的题面入口要先还原成具体任务：先统计频次，再只维护前 K 个候选。

\textbf{样例入口。} 拿 [1, 1, 1, 2, 2, 3]、k=2 手算，频次是 1->3、2->2、3->1，前两个高频是 1 和 2。

\textbf{暴力基线。} 暴力先统计频次，再把所有不同元素按频次完整排序，取前 K 个。

\textbf{暴力过程。} 以 [1, 1, 1, 2, 2, 3]、k=2 为例，频次为 1->3、2->2、3->1，答案是 1 和 2。

\textbf{重复工作。} 题面模型：先统计频次，再只维护前 K 个候选。暴力版的慢点要落到本题：浪费在完整排序所有不同元素。只要前 K 高频，低于当前第 K 高频边界的元素可以丢弃。后面的优化只消掉这类重复工作，不能改变答案集合。

\textbf{优化条件。} 本题的关键观察是：小顶堆大小超过 K 就弹出最低频；桶排序利用频次上界。这句话必须能翻译成可维护状态；如果题目条件改变后观察不再成立，就不能继续套原来的写法。

\textbf{相邻状态。} 规律是先用哈希计数，再用大小为 k 的小顶堆或桶按频次取 Top K；频次相同不影响题目集合答案。把这个特例继续拆开看：堆顶必须正好代表下一步最需要处理的候选；若堆里可能有过期对象，读取堆顶前要先清理。堆只解决动态极值，不自动证明被弹出或被丢弃的元素无用。

\textbf{状态复用。} 先用哈希表计数。再用大小为 K 的小顶堆保存候选，堆顶是当前前 K 中频次最低者；超过 K 就弹出。手算时只改被新输入影响的那一小部分，而不是回头重算完整候选。

\textbf{指针和字段。} 状态字段要能说清：freq[value] 是频次，heap 元素是 (frequency, value)，heap.top 是淘汰边界。手算时按这个协议跟踪：插入频次 3 和 2 后堆满；频次 1 的元素入堆会被立刻弹出，留下 1 和 2。

\textbf{代码落点。} 关键是比较器按频次建小顶堆。题目通常不要求输出顺序；若要求顺序，最后还要按规则排序结果。关键代码行必须能逐句翻译成“旧状态怎样变成新状态”，否则就只是模板。

\textbf{正确性。} 堆始终保存已处理元素中频次最高的 K 个；任何被弹出的元素频次不高于堆中 K 个候选，不可能进入答案。

\textbf{边界实验。} 先用K 等于不同元素数、频次相同、结果顺序不要求做反例检查。实验时覆盖K 等于不同元素数、频次相同、结果顺序不要求，记录堆顶是否过期、比较器是否让正确候选在堆顶、K 或窗口大小为边界值时是否稳定。

\textbf{复杂度变化。} 完整排序 O(m log m)，m 为不同元素数；堆法 O(n+m log k)，空间 O(m+k)。

\textbf{迁移追问。} 如果题目改成“若要按频次和字典序排序，比较器要同时处理两个键”，先判断原状态是否仍然覆盖新答案集合，再决定是微调边界、改 value 语义，还是换算法。

\subsection{LeetCode 1122 数组的相对排序}

\textbf{题目说明。} LeetCode 1122 数组的相对排序 的题面入口要先还原成具体任务：arr2 给出部分值的自定义顺序，剩余值按自然升序。

\textbf{样例入口。} 拿 arr1=[2, 3, 1, 3, 2, 4, 6, 7, 9, 2, 19]、arr2=[2, 1, 4, 3, 9, 6] 手算，先按 arr2 顺序输出 2 的三个、1 的一个、4 的一个、3 的两个、9 和 6，剩余 7、19 升序跟在后面。

\textbf{暴力基线。} 暴力可以用自定义比较器排序 arr1：arr2 中的值按 arr2 位置排序，其他值按自然升序。

\textbf{暴力过程。} 以 arr1=[2, 3, 1, 3, 2, 4, 6, 7, 9, 2, 19]、arr2=[2, 1, 4, 3, 9, 6] 为例，先输出所有 2，再 1、4、3、9、6，剩余 7、19 升序。

\textbf{重复工作。} 题面模型：arr2 给出部分值的自定义顺序，剩余值按自然升序。暴力版的慢点要落到本题：浪费在比较排序。题目值域有限时，计数数组直接保存每个值出现次数，输出时按目标顺序消耗即可。后面的优化只消掉这类重复工作，不能改变答案集合。

\textbf{优化条件。} 本题的关键观察是：值域有限时用计数数组，先按 arr2 顺序输出计数，再按数值升序输出未出现于 arr2 的数。这句话必须能翻译成可维护状态；如果题目条件改变后观察不再成立，就不能继续套原来的写法。

\textbf{相邻状态。} 规律是值域有限时用计数表保存频次，输出顺序由 arr2 和自然升序两段共同决定。把这个特例继续拆开看：先标出已经确认的前缀或后缀，再标出未知区；能被丢掉的候选，必须是以后再也不会改变已确认区域语义的候选。这样才算从例子推出数组不变量，而不是只记一次操作。

\textbf{状态复用。} count[x] 记录 arr1 中 x 的频次。先遍历 arr2，把 count[x] 个 x 写入答案；再按数值从小到大输出剩余 count。手算时只改被新输入影响的那一小部分，而不是回头重算完整候选。

\textbf{指针和字段。} 状态字段要能说清：count 是频次表，write 是答案写入位置，arr2 决定第一段自定义顺序。手算时按这个协议跟踪：count[2]=3，所以第一段先写三个 2；arr2 消耗完后，count[7] 和 count[19] 仍未归零，按自然升序写到末尾。

\textbf{代码落点。} 关键是输出 arr2 中值后把 count[value] 清零，避免剩余阶段重复输出。关键代码行必须能逐句翻译成“旧状态怎样变成新状态”，否则就只是模板。

\textbf{正确性。} 计数表完整保存 arr1 多重集合；第一段严格按 arr2 顺序输出，第二段按值递增输出不在 arr2 中的剩余值，正好符合题意。

\textbf{边界实验。} 先用arr2 元素互不相同、arr1 中有不在 arr2 的值、重复值、值域边界、空数组做反例检查。实验时把arr2 元素互不相同、arr1 中有不在 arr2 的值、重复值、值域边界、空数组列成表，再打印关键下标和有效区间。若某个元素被写入后又被未来步骤破坏，说明区间不变量没有成立。

\textbf{复杂度变化。} 排序 O(n log n)；计数法 O(n+值域+arr2长度)，空间 O(值域)。

\textbf{迁移追问。} 如果题目改成“若值域很大，改用哈希计数加排序剩余值或自定义比较器”，先判断原状态是否仍然覆盖新答案集合，再决定是微调边界、改 value 语义，还是换算法。

\subsection{LeetCode 283 移动零}

\textbf{题目说明。} LeetCode 283 移动零给定数组 nums，要求原地把所有 0 移到末尾，同时保持非零元素的相对顺序。

\textbf{样例入口。} 样例 nums=[0, 1, 0, 3, 12]，处理后为 [1, 3, 12, 0, 0]；非零元素 1、3、12 的相对顺序不变。

\textbf{暴力基线。} 暴力可以新建数组，先放所有非零元素，再补同样数量的零，最后复制回原数组。

\textbf{暴力过程。} 以 [0, 1, 0, 3, 12] 为例，辅助数组先收集 [1, 3, 12]，再补 [0, 0]。

\textbf{重复工作。} 题面模型：给定数组，原地把所有 0 移到末尾，同时保持非零元素的相对顺序。暴力版的慢点要落到本题：浪费在辅助数组。非零元素的相对顺序由读取顺序天然保证，只需要一个原地写入位置。后面的优化只消掉这类重复工作，不能改变答案集合。

\textbf{优化条件。} 本题的关键观察是：稳定保留非零元素顺序，把零集中到尾部；慢指针写入下一个非零位置，最后补零或用交换。这句话必须能翻译成可维护状态；如果题目条件改变后观察不再成立，就不能继续套原来的写法。

\textbf{相邻状态。} 规律是先稳定压缩非零前缀，再填零；不能在扫描时盲目交换导致顺序变化。把这个特例继续拆开看：先标出已经确认的前缀或后缀，再标出未知区；能被丢掉的候选，必须是以后再也不会改变已确认区域语义的候选。这样才算从例子推出数组不变量，而不是只记一次操作。

\textbf{状态复用。} write 指向下一个非零元素应写入的位置。第一遍把非零元素稳定写到前面，第二遍把 write 之后填 0。手算时只改被新输入影响的那一小部分，而不是回头重算完整候选。

\textbf{指针和字段。} 状态字段要能说清：read 是扫描位置，write 左侧是不含 0 的稳定前缀，write 到末尾最后会被置零。手算时按这个协议跟踪：读到 1 写到 nums[0]，读到 3 写到 nums[1]，读到 12 写到 nums[2]，最后 nums[3]、nums[4] 置 0。

\textbf{代码落点。} 关键代码是 if (nums[read]!=0) nums[write++]=nums[read]; 后面 while(write<n) nums[write++]=0。关键代码行必须能逐句翻译成“旧状态怎样变成新状态”，否则就只是模板。

\textbf{正确性。} 所有非零元素按原读取顺序写入前缀，因此稳定性保持；剩余位置数量等于原零的数量，补零后满足题意。

\textbf{边界实验。} 先用全零、无零、零在开头结尾、稳定顺序要求做反例检查。实验时把全零、无零、零在开头结尾、稳定顺序要求列成表，再打印关键下标和有效区间。若某个元素被写入后又被未来步骤破坏，说明区间不变量没有成立。

\textbf{复杂度变化。} 辅助数组 O(n) 空间；原地稳定写 O(n) 时间、O(1) 额外空间。

\textbf{迁移追问。} 如果题目改成“若目标是移动指定值，模型与移除元素相同”，先判断原状态是否仍然覆盖新答案集合，再决定是微调边界、改 value 语义，还是换算法。

\subsection{LeetCode 452 用最少数量的箭引爆气球}

\textbf{题目说明。} LeetCode 452 用最少数量的箭引爆气球 的题面入口要先还原成具体任务：一支箭的位置可以覆盖所有包含该位置的区间。

\textbf{样例入口。} 拿 [[10, 16], [2, 8], [1, 6], [7, 12]] 手算，按右端排序后第一箭射在 6，可以覆盖 [1, 6] 和 [2, 8]；下一箭射在 12。

\textbf{暴力基线。} 暴力可以枚举箭的位置，或对每个气球单独射箭后再尝试合并覆盖关系。

\textbf{暴力过程。} 以 [[10, 16], [2, 8], [1, 6], [7, 12]] 为例，按右端排序后第一支箭射在 6，可同时覆盖 [1, 6] 和 [2, 8]。

\textbf{重复工作。} 题面模型：一支箭的位置可以覆盖所有包含该位置的区间。暴力版的慢点要落到本题：浪费在没有利用区间右端。第一支箭若要覆盖最早结束的气球，射在它右端能给后续重叠区间最大机会。后面的优化只消掉这类重复工作，不能改变答案集合。

\textbf{优化条件。} 本题的关键观察是：按右端点排序，当前箭放在最早结束气球的右端。这句话必须能翻译成可维护状态；如果题目条件改变后观察不再成立，就不能继续套原来的写法。

\textbf{相邻状态。} 规律是每次把箭放在最早结束气球的右端，能最大化后续空间。把这个特例继续拆开看：先构造一个非贪心选择，再比较两者留下的资源、右边界或剩余时间。只有能把非贪心第一步交换成贪心第一步且不变差，局部选择才是规律。

\textbf{状态复用。} 按 right 升序排序。arrow\_pos 初始为第一个区间右端；后续区间若 start<=arrow\_pos 就被当前箭覆盖，否则新增一支箭。手算时只改被新输入影响的那一小部分，而不是回头重算完整候选。

\textbf{指针和字段。} 状态字段要能说清：arrow\_pos 是当前箭的位置，arrows 是箭数，interval 是当前气球闭区间。手算时按这个协议跟踪：第一箭在 6，覆盖 [1, 6] 和 [2, 8]；遇到 [7, 12] 时 7>6，需要第二箭放到 12。

\textbf{代码落点。} 关键是闭区间边界用 start<=arrow\_pos 表示能覆盖。坐标可能为负，不影响排序和比较。关键代码行必须能逐句翻译成“旧状态怎样变成新状态”，否则就只是模板。

\textbf{正确性。} 最早结束区间必须被某支箭覆盖；把这支箭移动到它的右端不会减少已覆盖区间，还可能覆盖更多后续区间。

\textbf{边界实验。} 先用端点相同、负坐标、区间包含、闭区间边界做反例检查。实验时除了端点相同、负坐标、区间包含、闭区间边界，还要故意替换成一个看似合理但错误的局部规则，构造反例。能解释错误贪心为何失败，才算理解正确贪心。

\textbf{复杂度变化。} 暴力选择箭位置复杂；贪心排序 O(n log n)，扫描 O(n)，空间取决于排序。

\textbf{迁移追问。} 如果题目改成“它和无重叠区间是同一类最早结束贪心”，先判断原状态是否仍然覆盖新答案集合，再决定是微调边界、改 value 语义，还是换算法。

\subsection{LeetCode 986 区间列表的交集}

\textbf{题目说明。} LeetCode 986 区间列表的交集 的题面入口要先还原成具体任务：两个有序且不重叠的区间列表，交集只可能来自当前两个区间。

\textbf{样例入口。} 拿 A=[0, 2], [5, 10] 和 B=[1, 5], [8, 12] 手算，当前两个区间交集是 [1, 2]，然后推进右端较小的 [0, 2]。

\textbf{暴力基线。} 慢版本对新区间和所有旧区间两两比较，手工判断相交、覆盖或冲突。它能验证端点语义，但排序后很多远离当前边界的区间无需再看。放到本题就是：先按“两个有序且不重叠的区间列表，交集只可能来自当前两个区间”完整试慢版本；样例里已经能看到“规律是两个有序列表只需要比较当前区间，谁先结束谁不可能再和对方后续产生交集”，所以优化前必须先能说清慢版本枚举了哪些候选。

\textbf{暴力过程。} 拿 A=[0, 2], [5, 10] 和 B=[1, 5], [8, 12] 手算，当前两个区间交集是 [1, 2]，然后推进右端较小的 [0, 2]。

\textbf{重复工作。} 题面模型：两个有序且不重叠的区间列表，交集只可能来自当前两个区间。暴力版的慢点要落到本题：慢版本会围绕“两个有序且不重叠的区间列表，交集只可能来自当前两个区间”做两两区间比较。样例规律是“规律是两个有序列表只需要比较当前区间，谁先结束谁不可能再和对方后续产生交集”，说明排序后只有贴近当前端点、结果尾部或资源堆顶的区间还会影响答案。后面的优化只消掉这类重复工作，不能改变答案集合。

\textbf{优化条件。} 本题的关键观察是：每次比较两个当前区间的重叠部分，然后推进右端较小的区间。这句话必须能翻译成可维护状态；如果题目条件改变后观察不再成立，就不能继续套原来的写法。

\textbf{相邻状态。} 规律是两个有序列表只需要比较当前区间，谁先结束谁不可能再和对方后续产生交集。把这个特例继续拆开看：排序以后，真正还能影响当前区间的候选必须贴近当前端点、结果尾部或资源堆顶。某个区间一旦被覆盖、合并或弹出，就要说明它为什么不会和后续区间重新产生更优关系。

\textbf{状态复用。} 把两两比较压成排序后的相邻关系、当前边界或动态有序结构；远离当前边界的区间要被安全归档。本题落点是：规律是两个有序列表只需要比较当前区间，谁先结束谁不可能再和对方后续产生交集。手算时只改被新输入影响的那一小部分，而不是回头重算完整候选。

\textbf{指针和字段。} 状态字段要能说清：排序后的区间、当前结果尾、扫描右端或资源堆、端点比较规则；动态版本还要保存前驱后继。手算时按这个协议跟踪：拿 A=[0, 2], [5, 10] 和 B=[1, 5], [8, 12] 手算，当前两个区间交集是 [1, 2]，然后推进右端较小的 [0, 2]。然后按协议复查：手算时先排序，再逐个区间写当前结果尾部、当前右端或资源堆。每个区间离开视野时都要说明它不会再影响后续。

\textbf{代码落点。} 排序比较器满足严格弱序；端点相等的处理和题意一致；扫描时只和当前结果尾或当前资源集合交互。关键代码行必须能逐句翻译成“旧状态怎样变成新状态”，否则就只是模板。

\textbf{正确性。} 证明从排序后的邻接关系开始：已经合并、选择或丢弃的区间不会被更后面的区间重新影响。

\textbf{边界实验。} 先用空列表、端点相接、一个区间覆盖多个区间、无交集做反例检查。实验时覆盖空列表、端点相接、一个区间覆盖多个区间、无交集，画出排序后的区间序列和当前结果尾部。动态版本要专门测前驱、后继和端点相接。

\textbf{复杂度变化。} 暴力两两比较区间常见为平方级；排序后扫描通常由排序成本主导，动态有序结构则把每次前驱后继查询控制在对数级。

\textbf{迁移追问。} 如果题目改成“这是区间双指针，不需要把两个列表合并排序”，先判断原状态是否仍然覆盖新答案集合，再决定是微调边界、改 value 语义，还是换算法。

\begin{principlebox}{本节复盘方式}
不要只看最终算法名。每道题复盘时先遮住优化版，口述暴力枚举对象；再指出相邻窗口、历史前缀、候选区间、搜索状态或子问题里哪一部分被重复处理；最后用一个边界样例验证状态更新顺序。能讲完这三步，才算真正掌握本章案例。
\end{principlebox}
